\documentclass[oneside,titlepage,a4paper]{report} %book,article,report,letter
\usepackage{amsmath} %宏包
\usepackage{graphicx}
\usepackage{amssymb}
\usepackage{titling}
\usepackage{hyperref}
\usepackage{pdfpages}
\usepackage{xcolor}
\usepackage{enumerate}
\usepackage{float}
\usepackage{multirow} 
\usepackage{shorttoc}
\usepackage{tabularx}

%\usepackage{geometry}

%\geometry{left=2.0cm,right=2.0cm,top=2.0cm,bottom=2.0cm}
%\usepackage[T1]{fontenc}
%\usepackage{makecell}
%\renewcommand\theadfont{\bfseries\sffamily}
%导言区

\setcounter{tocdepth}{4}
\setcounter{secnumdepth}{3}
%header and footer setting
\usepackage{fancyhdr}
\pagestyle{fancy}
\fancyhf{}
\rhead{Research on GameSpy protocol}
\lhead{\includegraphics[scale=0.05]{Pictures/RetroSpyLogoWithText.png}}
\lfoot{Page \thepage}
\rfoot{\includegraphics[scale=0.07]{Pictures/RetroSpy_Logo_with_Long_Text.png}}
\fancypagestyle{plain}{\pagestyle{fancy}}

\newcommand{\myparagraph}[1]{\paragraph{#1}\mbox{}\\}

\begin{document}

\title{\Huge\textbf{Research On GameSpy Protocol}} 
\author{Arves100, Xiaojiuwo}

\renewcommand\maketitlehooka{
	\pagestyle{empty}
	\centering
	}
\renewcommand\maketitlehookc{	
	\includegraphics[width=5cm]{Pictures/GameSpylogoindustries.pdf}
	\vfill First Edition}


%\date{} %%如果没有这句，会生成时间

\maketitle  %%生成书名

\tableofcontents  %%生成目录

%\mainmatter %%表示文章的正文部分，在生成目录后将从第一页开始

\part{Research On GameSpySDK}

\chapter{GameSpy General Construction}
\par In GameSpy SDK there are 16 modules, which constructed the GameSpy main functions.
\section{GameSpy SDK Module}
\begin{itemize}
	\item GameSpy Presence Servers
	\begin{itemize}
		\item GameSpy Presence Connection Manager
		\item GameSpy Presence Search Player
	\end{itemize}
	\item Nat Negotiation
	\item Master Server: Query Report 2
	\item Master Server: Server Browser
	\item Master Server: Available Check
	\item Game Patching
	\item Game Tracking
	\item Master Server Patching: Downloading files from FilePlanet
	\item Peer SDK
	\item Game Statitics
	\item Chat Server
\end{itemize}

\chapter{GameSpy Presence Servers}
\par GameSpy Presence Servers contain two server, GameSpy Presence Connection Manager (GPCM) and GameSpy Presence Search Player (GPSP). GPCM is a server that handle login request and response with corresponding user infomation stored on GameSpy. GPSP is a server that handle search request for user.
\subsection{Server IP and Ports}
Table \ref{IP and Ports for GameSpy Presence Servers} are the GPCM and GPSP IP and Ports that client/game connect to.
\begin{table}[H]
	\centering
	\begin{tabular}{|c|c|}
		\hline 
		IP&Port  \\ 
		\hline 
		gpcm.gamespy.com&29900 \\ 		
		\hline 
	 	gpsp.gamespy.com&29901 \\
	 	\hline 
	\end{tabular} 
\caption{IP and Ports for GameSpy Presence Servers}
\label{IP and Ports for GameSpy Presence Servers}

\end{table}

\subsection{Database Key Field}
These keys is that GameSpy Presence SDK using to find a user in their database. Keys are shown in Table \ref{Key Field}.

\begin{table}[H]
	\centering
	\begin{tabular}{|c|>{\centering\arraybackslash}p{8cm}|}
		\hline 
		Keys& Description  \\ 
		\hline 
		User & An user contains the Email and the password, but contains multiple profiles \\ 		
		\hline 
		ProfileID & The profile contains the name, surname, birth date and all the rest user info, including
		an unique nickname used to identify the profile and a generic nickname used to show for example in
		games \\
		\hline 
	\end{tabular} 
	\caption{Key Field}
	\label{Key Field}	
\end{table}

\subsection{Protocol Descriptions}

In this part, we show the protocol detail in GameSpy Presence SDK.
\subsubsection{The String Pattern}
We first introduce the pattern of the string, which is using to make up a request.
This kind of string is represent a value in a request sends by the client as Table \ref{Value string}.

\begin{table}[H]
	\centering
	\begin{tabular}{|c|c|}
		\hline 
		String&Description  \\ 
		\hline 
	$ \backslash \langle content \rangle \backslash $& The value is $ \langle content \rangle $  \\ 
		\hline 
	\end{tabular} 
\caption{Value string}
\label{Value string}
\end{table}

This kind of string is represent a command in a request sends by the client as Table \ref{Command string}. The command will end with $ \backslash \backslash $ or $ \backslash $ depends on whether run at the server-side or client-side.


\begin{table}[H]
	\centering
	\begin{tabular}{|c|c|}
		\hline 
		String&Description  \\ 
		\hline 
		$ \backslash command \backslash\backslash $& This is a command \\ 		
		\hline 
		$ \backslash error \backslash \backslash $ & Error command \\
		 \hline
		 $\backslash lc \backslash$& Login command\\
		 \hline
	\end{tabular} 
	\caption{Command string}
	\label{Command string}
\end{table}

This kind of string is represent a parameter in a request sends by the client \ref{Parameter string}. GameSpy uses the combination of the parameter to search the string with value, and sends the data back to client use this kind of parameter string.

\begin{table}[H]
	\centering
	\begin{tabular}{|c|c|}
		\hline 
		String&Description  \\ 
		\hline 
		$ \backslash id \backslash 1 \backslash $& This is a parameter string the value of $ id $ is $ 1 $ \\ 		
		\hline 
		$ \backslash profileid \backslash 007 \backslash $ & This is a parameter string the value of $ profileid $ is $ 007 $ \\
		\hline
	\end{tabular} 
	\caption{Parameter string}
	\label{Parameter string}
\end{table}

Error response string for (GPCM, GPSP):
\begin{equation}
\begin{split}
\backslash error \backslash\backslash err \backslash \langle error code \rangle \backslash fatal\backslash\backslash errmsg \backslash \langle error message \rangle \backslash id\backslash 1 \backslash final \backslash
\end{split}	
\end{equation}
\subsubsection{Login Phase}
\myparagraph{Client Login Request}
There are three ways of login:
\begin{itemize}
	\item AuthToken: Logging using an alphanumeric string that rapresents an user
	\item 	UniqueNick: Logging using a nickname that is unique from all the players
	\item User: Logging with the nickname and the password
\end{itemize}

The full login request string:
\begin{equation}\label{Chanllenge string}
\begin{gathered}
	\backslash login \backslash challenge \backslash \langle challenge \rangle \backslash authtoken \backslash \langle authtoken \rangle \\ \backslash uniquenick \backslash \langle uniquenick \rangle \backslash user \backslash \langle user \rangle 
	\backslash userid \backslash \langle userid \rangle \\ \backslash profileid \backslash \langle profileid \rangle \backslash partnerid \backslash \langle partnerid \rangle \backslash response \backslash \langle response \rangle \\
	 \backslash firewall \backslash 1 \backslash port \backslash \langle port \rangle \backslash productid \backslash  \langle productid \rangle \\ \backslash gamename \backslash \langle gamename \rangle \backslash namespaceid \backslash \langle namespaceid \rangle \\ \backslash  sdkrevision \backslash \langle sdkrevision \rangle \backslash quiet \backslash \langle quiet \rangle \backslash id \backslash 1 \backslash final \backslash
\end{gathered}
\end{equation}
The value $ \langle challenge \rangle $ for $ \backslash challenge \backslash $ in \ref{Chanllenge string} is a 10 byte alphanumeric string.

The following Table \ref{Login parameter string} is a description of string used in login request, GameSpy can use these string to find value in database.
\begin{table}[H]
	\centering
	\begin{tabular}{|c|>{\centering\arraybackslash}p{9cm}|}
		\hline
		  String    & Description                                                                                                                                                                                               \\ \hline
		 login& The login command which use to identify the login request of client\\ \hline
		 challenge  & The user challenge used to verify the authenticity of the client                                                                                                                                          \\ \hline
		 authtoken  & The token used to login (represent of an user)                                                                                                                                                            \\ \hline
		uniquenick  & The unique nickname used to login                                                                                                                                                                         \\ \hline
		   user     & The users account (format is NICKNAME@EMAIL)                                                                                                                                                              \\ \hline
		  userid    & Send the userid (for example when you disconnect you will keep this)                                                                                                                                      \\ \hline
		 profileid  & Send the profileid (for example when you disconnect you will keep this)                                                                                                                                   \\ \hline
		 partnerid  & This ID is used to identify a backend service logged with gamespy.(Nintendo WIFI Connection will identify his partner as 11, which means that for gamespy, you are logging from a third party connection) \\ \hline
		 response   & The client challenge used to verify the authenticity of the client                                                                                                                                        \\ \hline
		 firewall   & If this option is set to 1, then you are connecting under a firewall/limited connection \\
		 \hline
		 port& The peer port (used for p2p stuff)                                                                          \\ \hline
		 productid  & An ID that identify the game you're using                                                                                                                                                                 \\ \hline
		 gamename   & A string that rapresents the game that you're using, used also for several activities like peerchat server identification                                                                                 \\ \hline
		
		namespaceid & ?                                                                                                                                                                                                         \\ \hline
		sdkrevision & The version of the SDK you're using \\ \hline
		   quiet    & ? Maybe indicate invisible login which can not been seen at friends list \\ \hline
		   id& The value is 1\\ \hline
		   final & Message end\\ \hline
	\end{tabular} 
	\caption{Login parameter string}
	\label{Login parameter string}
\end{table}
\paragraph{Login Response From Server}\mbox{}\\
There are two kinds of login response string: 
\begin{equation}\label{Login response string1}
\begin{split}
\backslash lc \backslash 1
\end{split}	
\end{equation}
\begin{equation}\label{Login response string2}
\begin{split}
\backslash lc \backslash 2
\end{split}	
\end{equation}


This response string \ref{Login response string1}, \ref{Login response string2} is send by the server when a connection is accepted, and followed by a challenge\ref{Chanllenge string}, which verifies the server that client connect to.



\subsubsection{User Creation}
This commmand \ref{Create user command} is used to create a user in GameSpy.
\begin{equation}\label{Create user command}
\begin{gathered}
\backslash newuser \backslash email \backslash \langle email \rangle \backslash nick \backslash \langle nick \rangle \backslash passwordenc \backslash \langle passwordenc \rangle \\
\backslash productid \backslash \langle productid \rangle \backslash gamename \backslash \langle gamename \rangle \backslash uniquenick \backslash \langle uniquenick \rangle \\ \backslash cdkeyenc \backslash \langle cdkeyenc \rangle \backslash partnerid \backslash \langle partnerid \rangle \backslash id \backslash 1 \backslash final \backslash
\end{gathered}	
\end{equation}


\part{RetroSpy System Architecture}


\chapter{introduction}

\chapter{conclusion}

\end{document}
